\documentclass[11pt,a4paper]{article}

\usepackage[margin=1.05in]{geometry}
\usepackage{amsmath,amssymb,amsthm}
\usepackage{mathtools}
\usepackage{stmaryrd}
\usepackage{microtype}
\setlength{\emergencystretch}{2.5em}
\usepackage{enumitem}
\usepackage{booktabs}
\usepackage{array}
\usepackage{xcolor}
\usepackage[colorlinks=true,linkcolor=blue!45!black,urlcolor=blue!60!black]{hyperref}

\newtheorem{theorem}{Theorem}[section]
\newtheorem{proposition}[theorem]{Proposition}
\newtheorem{corollary}[theorem]{Corollary}
\newtheorem{principle}[theorem]{Principle}
\theoremstyle{definition}
\newtheorem{definition}[theorem]{Definition}
\newtheorem{example}[theorem]{Example}
\theoremstyle{remark}
\newtheorem{remark}[theorem]{Remark}
\newtheorem{reading}[theorem]{Reading}
\newtheorem{question}[theorem]{Question}

\newcommand{\FF}{\mathbb{F}}
\newcommand{\NN}{\mathbb{N}}
\newcommand{\ZZ}{\mathbb{Z}}
\newcommand{\cat}[1]{\mathbf{#1}}
\newcommand{\Set}{\cat{Set}}
\newcommand{\Vect}{\cat{Vect}}
\newcommand{\ext}[1]{\llbracket #1 \rrbracket}
\newcommand{\con}{\triangleleft}
\DeclareMathOperator{\Hom}{Hom}
\DeclareMathOperator{\Cov}{Cov}
\newcommand{\Schub}{\mathfrak{S}}

\definecolor{okgreen}{rgb}{0.0,0.45,0.15}
\definecolor{nored}{rgb}{0.65,0.0,0.0}
\newcommand{\OK}{\textcolor{okgreen}{\textbf{confirmed}}}
\newcommand{\NO}{\textcolor{nored}{\textbf{refuted}}}

\title{\bfseries Two Planes\\[0.4em]
\large the data/control distinction, and the three things it forces}
\author{Robin Langer \and Claude}
\date{21 September 2026}

\begin{document}
\maketitle

\begin{abstract}
A container morphism is a pair $(u, f)$ of opposite variance: $u$ dispatches on
shapes, $f$ routes positions, and $f$ runs backwards. Shapes form a bare set and
positions carry structure, so the two components live in different categories.
This note argues that the asymmetry is not incidental notation but the load-bearing
wall of the formalism, and that one prohibition follows from it --- \emph{a
container morphism may not look at the payload} --- from which three otherwise
unrelated results follow.

First, redundancy is available on one plane only: coding needs an addition, and
shapes have none, so the control plane must be protected by replication rather
than by codes (\texttt{error-correcting.pdf}). Second, a federation carries two
independent obstructions, not one --- well-foundedness in $u$ and a class in
$H^1$ in $f$ --- which resolves a question left open in
\texttt{sheaf-federation.pdf} and is verified on running servers. Third, the
flag container's Hecke algebra is an algebra of \emph{dispatch}, which is why
Schubert polynomials, needing graded positions, do not appear
(\texttt{schubert.pdf}).

We then record what cannot be repaired: linearising the shapes collapses the two
planes into one and forgets the shape set entirely. And we state the limit
honestly --- real systems need data to influence control, so the formalism
describes not a whole system but the fragment that can be checked statically. Its
value is that it tells you exactly where that fragment ends.
\end{abstract}

\section{The distinction}

\begin{definition}[the two planes]
A container is $S \con P$: a set $S$ of \emph{shapes}, and for each $s$ an object
$P s$ of \emph{positions}. A morphism $(S \con P) \to (T \con Q)$ is a pair
\[
  u : S \longrightarrow T,
  \qquad
  f_s : Q(u\,s) \longrightarrow P s .
\]
We call $u$ the \emph{control plane} and $f$ the \emph{data plane}.
\end{definition}

Three features, and each does work below.

\begin{enumerate}[nosep,label=(\roman*)]
\item \textbf{Opposite variance.} Control flows forward, data flows back.
\item \textbf{Different categories.} $S$ lives in $\Set$; $P s$ lives in $\Vect$,
      in $\FF_q$-spaces, or in whatever the base is.
\item \textbf{One-way dependence.} $f$ may depend on $s$. $u$ may not depend on
      the data.
\end{enumerate}

\begin{principle}[the prohibition]\label{prin:prohibit}
A container morphism may not look at the payload. Equivalently: there is no
data $\to$ control flow.
\end{principle}

Everything in this note is a consequence of Principle~\ref{prin:prohibit} or of
the fact that (ii) is what sustains it.

\section{First instance: redundancy lives on one plane}

\begin{theorem}[\texttt{error-correcting.pdf}]\label{thm:code}
Linear redundancy is available for positions and unavailable for shapes.
\end{theorem}

\begin{proof}[Reason]
Redundancy means forming linear combinations, which requires an addition on the
values protected. Positions span a vector space and have one. Shapes form a bare
set: there is no $s_1 + s_2$.
\end{proof}

\begin{principle}[two planes, two mechanisms]
Protect positions with codes. Protect shapes with replication, acknowledgement or
consensus.
\end{principle}

The check that this is more than bookkeeping is that practice already obeys it and
did not know why. Packet headers are protected by checksum-and-discard, never by
coding across headers --- a corrupted header means the packet is dropped, not
repaired. Storage systems replicate metadata and erasure-code data. A system that
erasure-codes its payload and leaves its metadata on a single unreplicated path
has not been made reliable; it has moved its single point of failure into the
header. The formalism says this was not a design choice.

\section{Second instance: a federation has two obstructions}

This is the new claim of the present note, and it reinterprets a result obtained
independently in \texttt{trace-and-class.pdf}.

\emph{A Sheaf Theory of Federation} asked whether the trace obstruction on
$\con$ and the class in $H^1$ are the same obstruction, offering two answers:
they are the same, or one of the two models is wrong. Both are too tight.

\begin{proposition}\label{prop:fed}
In the ring of container morphisms, well-foundedness is a property of $u$ and the
cohomology class is a property of $f$. Since the two components of a morphism are
chosen independently, the two obstructions are independent.
\end{proposition}

The evidence is that the four handlers of the reference implementation all apply
the transition function $\varepsilon$ in the same position --- to the reply, on
the way back, which is the only position descent allows --- and differ only in
what stops them: fuel, fuel recorded in the response type, a visited list, or
nothing at all. What distinguishes the diverging handler from the terminating
ones is therefore a property of the forward map's control flow, where no
$\varepsilon$ occurs. Writing the recursion out,
$\mathrm{Orient}(i) = \varepsilon_i \cdot \mathrm{Orient}(i+1)$, the argument
never mentions $\varepsilon$.

\begin{center}
\small
\begin{tabular}{llcccl}
\toprule
class & repair & $H^1$ & terminates & canonical & witness \\
\midrule
cylinder & none & $0$      & no  & ---  & abandoned at 8s, 757 hops \\
cylinder & cut  & $0$      & yes & yes  & all three cuts agree \\
mobius   & none & $\ne 0$  & no  & ---  & abandoned at 6s, 660 hops \\
mobius   & cut  & $\ne 0$  & yes & no   & the three cuts disagree \\
\bottomrule
\end{tabular}
\end{center}

All four cells populate, on three live servers. The load-bearing cell is the
first: on the cylinder the twisted sheaf has two global sections, nothing
whatever obstructs descent, and the unbounded walk still does not return. A
vanishing class did not buy termination, because termination was never what the
class was about.

\begin{corollary}
The two planes obstruct different properties. Well-foundedness governs
\emph{existence} --- is there an answer at all. The class governs
\emph{canonicity} --- is the answer independent of the choice made to obtain it.
Cutting an edge trades divergence for choice, and $H^1$ is exactly the price of
the choice.
\end{corollary}

Checked by enumeration on all eight elements of $(\ZZ/2)^3$, together with the
gauge invariance --- all $24$ gauge transformations move the cocycle and never
move the verdict --- which is what makes the statement one about $H^1$ rather
than about a chosen representative.

\begin{remark}
The sheaf model sees only the data plane, because $H^1$ is all it offers. The
operational model sees only the control plane, because divergence is all it
offers. Neither is wrong; each is complete on its own plane. That is why the
question as posed had no good answer.
\end{remark}

\section{Third instance: the flag container's algebra is an algebra of dispatch}

The flag container is $\sum_{w \in S_n} X^{\ell(w)}$: shapes are the permutations,
positions over shape $w$ are the Bruhat cell $\FF_q^{\ell(w)}$. So the Weyl group,
the Bruhat order and the Hecke algebra all live on the \emph{control} plane, and
the cells are the data.

\texttt{schubert.pdf} shows that the fibration forgetting the $i$-th subspace
gives $S_i$ with $S_i^2 = (q+1)S_i$, hence $T_i = S_i - I$ satisfying the Hecke
quadratic, with the braid relations holding and no nilpotent attached to any
single index. Read through the present distinction, that result says: \emph{the
algebra the container produces is the algebra of how it dispatches.}

\begin{corollary}
Schubert polynomials are not reachable because they are a data-plane refinement.
$\Schub_w(x_1, \ldots, x_n)$ records the torus weights on positions; counting
$\FF_q$-points records only how many positions there are. The one-variable shadow
of $\Schub_w$ is $x^{\ell(w)}$, and the whole container decategorifies to the
Poincar\'e polynomial $\sum_w x^{\ell(w)}$, which at $x = q$ is $[n]_q!$.
\end{corollary}

\begin{question}
Grade the positions by the torus action, so that a position carries a weight and
not merely a dimension. Does the multigraded flag container have
$\Schub_w(x_1,\ldots,x_n)$ as its shape polynomial?
\end{question}

\section{Why the distinction cannot be repaired away}

It is tempting to regard the set-ness of shapes as a limitation to be generalised
away by making both planes linear. \texttt{q-container.pdf} Theorem~B says what
happens.

\begin{theorem}[Theorem~B, restated for this purpose]
If the extension is taken to be $\bigoplus_s \Hom(P s, V)$ rather than
$\sum_s \Hom(P s, V)$, the natural transformations collapse to
$\Hom\bigl(\bigoplus_t Q t, \bigoplus_s P s\bigr)$, which forgets the shape set
entirely.
\end{theorem}

\begin{corollary}\label{cor:wall}
The distinction survives precisely because the two planes live in different
categories. Collapse them into one and one does not obtain a more general theory;
one obtains a single plane, and the control plane is the one that disappears.
\end{corollary}

So (ii) is the load-bearing wall, and Theorem~\ref{thm:code} is not an accident of
presentation: shapes have no addition because shapes are in $\Set$, and they are
in $\Set$ because that is the only thing keeping two planes from becoming one.

\section{The agentic reading}

\begin{reading}[stated as a reading, not a theorem]
Prompt injection is data becoming control: a payload read by the dispatcher.
Principle~\ref{prin:prohibit} says a container morphism cannot do this, since $u$
is a function of the shape alone. The vulnerability class is therefore not
mitigated within the formalism; it is \emph{unexpressible} in it.
\end{reading}

Two qualifications, because the claim is easy to overstate. Unexpressible in the
fragment is not the same as absent from the system: a real agent does read its
inputs and branch, and that part of it is outside the formalism by
\S\ref{sec:outside}. What the discipline buys is that the unsafe flow must be
written somewhere explicit, rather than arising by accident anywhere.

And it explains the status of labels carried through a system. A trust label that
is transported unchanged by $f$ and never consulted by $u$ is a
\emph{passenger}: it needs no type discipline of its own, precisely because it
never reaches the control plane.

\section{What is outside, and the limit}\label{sec:outside}

\begin{center}
\small
\begin{tabular}{lll}
\toprule
operation & morphism? & why \\
\midrule
encoding, syndrome extraction $s = Hy$ & yes & a matrix \\
erasure decoding, pattern known & yes & pattern is metadata, recovery a matrix \\
error correction, position unknown & \textbf{no} & must infer from the payload \\
retransmission, rate adaptation & \textbf{no} & reads state and branches \\
\bottomrule
\end{tabular}
\end{center}

The line falls in exactly one place, every time: an operation is outside when data
influences control. Note the split \emph{inside} decoding --- computing a syndrome
is linear and therefore a morphism, while mapping a syndrome to an error vector is
a search and therefore not. The formalism separates the algebra of decoding from
the search.

\begin{remark}[the honest limit]
Real systems need data $\to$ control flow: retransmission, rate adaptation, error
correction, any adaptive agent. The formalism therefore does not describe a whole
system. It describes the fragment that can be checked statically, and its value is
that it tells you exactly where that fragment ends. As
\texttt{error-correcting.pdf} says plainly, it yields no new codes; what it yields
is recoverability as a property of declared interfaces, computable by multiplying
matrices, rather than established per-code and per-deployment.
\end{remark}

\begin{principle}[what to do with a boundary]
When a system needs data to influence control, that is not a failure to
formalise. It is a boundary one should be able to point at. Containers make you
point at it.
\end{principle}

\section{Status}

\begin{center}
\small
\setlength{\tabcolsep}{4pt}
\begin{tabular}{cp{0.40\textwidth}ll}
\toprule
\# & claim & source & verdict \\
\midrule
$1$ & redundancy unavailable for shapes & \texttt{error-correcting} & proved \\
$2$ & the two federation obstructions are independent & \texttt{trace-and-class} & \OK{}, live + $8/8$ \\
$3$ & flag container's algebra is Hecke, not nilCoxeter & \texttt{schubert} & proved + $5/5$ \\
$4$ & linearising shapes collapses the planes & \texttt{q-container} Thm~B & proved \\
$5$ & Prop.~\ref{prop:fed}: (1)--(4) are one distinction & this note & synthesis, new \\
$6$ & injection is unexpressible in the fragment & this note & reading, not proved \\
\bottomrule
\end{tabular}
\end{center}

\noindent\textbf{Companions.} \texttt{q-container.pdf},
\texttt{error-correcting.pdf}, \texttt{schubert.pdf} in this directory;
\texttt{trace-and-class.pdf} and \texttt{independence.ts} in
\texttt{gold/federation}.

\end{document}
